\scanpage{119}
% archive foliation: 118
Notons que si $X_0=\mathbb P^1(\mathbb C)$ \deleted{il y a} le faisceau $H\backslash\mathcal P_{X_0}$ est le faisceau final, \deleted{\unsure{et que}} \deleted{\unsure{il y a}} il y a une conn. projective can. \inserted{$c$} sur $X_0$, savoir $c_t$ pour $t=\operatorname{id}_{X_0}$ ! C'est la seule conn. projective sur $X_0$, car $H^0(X_0,\underline\omega_{X_0}^{\otimes2})=0$ ($\underline\omega$ est de degré $-2$, donc $\underline\omega^{\otimes2}$ de degré $-4$ !).

Revenant au cas d'un $X$ général, on voit donc que
\[
c_t=t^*(c)=\text{image inverse par }t\text{ de la conn. projective can. }c
\]
(l'unique conn. proj.) sur $\mathbb P^1(\mathbb C)$.
\deleted{\unsure{c}}

\uline{Extension à la géom. alg.} Soit maintenant $X$ schéma form$^{\mathrm t}$. lisse de dim \inserted{diff} relative 1 sur un schéma $S$. Si \deleted{$t,s:$}
\[
t,s:X\longrightarrow\mathbb E^1_S
\]
sont étales \inserted{\unsure{topologiques}}, on définit comme plus haut
\[
G(t,s)\in\Gamma(X,\underbrace{(\Omega^1_{X/S})^{\otimes2}}_{\underline\omega^{\otimes2}}),
\]
et on vérifie bien sûr (1) et (2) par le $=$ calcul que dans le cas analytique, ce qui permet d'étendre $G(t,s)$ au cas de morphismes \uline{étales}
\[
t,s:X\longrightarrow\mathbb P^1_S
\]
\deleted{\unsure{topologiques}}, de façon à satisfaire encore \deleted{$t,s:X\longrightarrow$}

(posant $H_S=\underline{\operatorname{Aut}}_S(\mathbb P^1_S)$)
\[
G(gt,g's)=G(t,s)\qquad(g,g'\in\Gamma(S,H))
\]
et à définir ainsi, si $\mathcal P_{X/S}$ est le faisceau \inserted{(ét)} des morphismes étales $X\longrightarrow\mathbb P^1_S$, \deleted{\unsure{défini par}} \inserted{un morphisme can.}
\[
G:\mathcal P_{X/S}\times\mathcal P_{X/S}\longrightarrow\underline\omega_{X/S}^{\otimes2}.
\]
\margin{\unsure{au lieu de diff. principales parler plutôt de diff. uniques} jusqu'à l'ordre 3, i.e. \unsure{que} $\unsure{\operatorname{gr}(P^3_{X/S})}$ \unsure{$\simeq\operatorname{Sym}(\underline\omega_{X/S})$} $\unsure{\longrightarrow}$ \unsure{iso en degrés} $\le3$.}
